\section{Introduction}

Local neighborhoods, graphs, and related structural signals are established
ways to shape unsupervised representations \citep{xu2024sccdcg,li2025dcboost}.
They can provide useful inductive bias for sparse, high-dimensional observations,
but a structural signal is not itself a clustering objective. The operation that
selects a relation, applies a perturbation, updates a representation, and then
produces a partition can preserve, distort, or discard the information available
at the start.

We study the gap between \emph{structural opportunity} and
\emph{intervention utility}. A graph may contain a locally informative relation
while a learned selector chooses the wrong coordinates. A fixed-budget
intervention may damage cluster-defining features. An auxiliary objective may
compete with representation learning \citep{yan2026scmib}. Or a local geometric
improvement may fail to survive the global readout \citep{li2025dcboost}.
These stages are often collapsed into one score difference between a structural
variant and a baseline. Such a difference shows that a variant changed a score;
it does not identify whether the change came from the structural policy, the
generic intervention, the optimizer, or the readout.

The problem becomes sharper when evidence accumulates across exploratory
versions. Different versions may change preprocessing, intervention location,
training target, effective budget, stochastic schedules, or final readout.
Seeds and variants are then repeated measurements rather than independent data
sets. A credible synthesis must preserve protocol and dataset boundaries while
still exposing recurring patterns. It must also distinguish post-treatment
descriptors from pre-intervention properties and must not turn incomplete
computation into a negative result.

Our question is therefore: \textbf{why does apparently useful structural
information fail to translate reliably into unsupervised clustering utility?}
We use the localization chain

\begin{center}
\textbf{Opportunity $\rightarrow$ Selection $\rightarrow$ Intervention
$\rightarrow$ Representation $\rightarrow$ Readout}.
\end{center}

The chain asks whether an opportunity is measurable, whether topology changes
what is selected, whether a matched intervention is helpful, whether the
objective moves the representation compatibly, and whether local changes survive
the global readout. This is a study protocol and claim firewall, not a new
neural architecture.

\begin{figure}[t]
\centering
\includegraphics[width=0.98\linewidth]{V25_Figure2_mechanism_chain}
\caption{Frozen V25 mechanism-localization chain. The diagram is a protocol
schema for separating opportunity, policy, intervention, representation, and
readout; it is not a fitted architecture or a causal graph estimated from the
historical rows.}
\label{fig:chain}
\end{figure}

Our contributions are twofold. First, we provide an auditable observational
Failure Atlas for the V1--V22 history, with repeated-measurement units,
protocol/readout boundaries, and a separate V23/V24 boundary branch. Second, we
provide a matched V21 selection-policy decomposition that separates generic
intervention from topology-dependent selection and exposes conditional signs.
Feature, gradient, and local/global analyses localize possible failure stages;
they are not promoted to universal causal laws.
